
District-heating networks serving industrial demand are central to European
decarbonisation, integrating heat pumps, combined heat and power and thermal storage in
multi-source systems~\cite{Lund2014_4GDH}, and their operational optimisation relies
increasingly on mathematical programming~\cite{PereaMoreno2026}. District heating already
meets a substantial share of building heat demand in mature markets~\citep{Werner2017} and
remains integral to clean-energy-system designs~\citep{Lund2025_clean}. Large heat pumps in
particular are becoming central to its decarbonisation~\citep{David2017}, which raises the
economic stakes of the operational models that schedule them. How much of the network
such a model represents is a design choice, and it is made along two axes: the spatial
resolution at which the network appears, from a single-bus copperplate to a pipe-by-pipe
graph; and the physical phenomena admitted, from a lossless energy balance to a model with
pressure-coupled hydraulics, temperature propagation and transport delay.

The two extremes are well populated. Copperplate models aggregate all components and
demands onto a single bus, ignoring spatial structure and thermal losses entirely, a
limitation their authors acknowledge: \citet{Wirtz2020} note that ``the
validity of the global energy balance assumption depends on the network topology, size and
the operation strategy'', without quantifying the resulting error. Full thermo-hydraulic
models resolve the network pipe by pipe and capture the thermal losses that amount to
5--15\,\% of annual heat supply~\cite{Frederiksen2013}, along with pressure-dependent
pumping, temperature propagation and the delays that follow from finite fluid velocity. A
third complication sits across both axes: extended physics introduces nonlinear and
bilinear terms incompatible with standard mixed-integer linear programming, forcing a
choice between linearisation, which preserves tractability at the cost of approximation
error, and native nonlinear formulations, which preserve the physics at the cost of solve
time.

The working assumption in practice is that finer spatial resolution and richer physics
both buy proportionate accuracy, and that the computational cost is therefore justified in
proportion. No quantitative evidence establishes that proportionality, or identifies which
axis delivers the return. The closest evidence comes from a neighbouring field: in
electricity systems, coarse spatial resolution causes capacity-expansion models to miss
intra-regional congestion and understate system cost, and the spatial aggregation choice
dominates other modelling simplifications in scenario
comparisons~\cite{Frysztacki2021}. A parallel power-systems analysis shows spatial and
temporal resolution trading off against one another under a fixed computational
budget~\citep{FrewJacobson2016}, the same class of resolution-versus-physics trade this
paper prices for district heating. It is natural to expect the analogue in district
heating, and that expectation is worth stating precisely because this paper finds it
holds only in part. Once the losses are supplied by other means, the spatial structure itself
contributes little to dispatch cost, at least where all generation is co-located at a
single source; whether that holds under distributed generation is left as a
controlled open question. The distinction is
invisible to any comparison that changes both at once, which is the first of the two
limitations this study addresses.
